\section{Background and Related Work}\label{sec:related}

\subsection{Illicit-activity detection on the Bitcoin transaction graph}
Weber \emph{et al.}~\cite{weber2019} released the Elliptic data set, a graph of roughly 200 thousand Bitcoin transactions with 166 anonymised features, 234 thousand payment-flow edges and labels for about a quarter of the nodes, organised in 49 time steps. Their own comparison already contained the result that most later work has had to contend with: a Random Forest on node features outperformed a graph convolutional network, and every model degraded after a dark-market shutdown late in the timeline. Subsequent studies applied temporal graph networks~\cite{pareja2020}, supervised learners~\cite{alarab2020}, gradient boosting~\cite{vassallo2021} and active learning under label scarcity~\cite{lorenz2020} to the same data; broader surveys place this line of work in the context of anti-money-laundering analytics~\cite{chen2018} and cryptocurrency network analysis~\cite{wu2021}.

Elmougy and Liu~\cite{elmougy2023} extended the data set to Elliptic++ by adding 17 transaction features and, more importantly, an \emph{actor} view: 822,942 wallet addresses with behavioural features, together with address-to-address and address-to-transaction edge lists. They report illicit-class F1 of 0.828 for transactions and 0.845 for wallets with a 50-tree Random Forest under a time-step split (steps 1--34 for training, 35--49 for testing). These figures are the natural reference point for our work and we return to them in Section~\ref{sec:res-published}.

\subsection{Does graph structure help?}
Message-passing networks such as GCN~\cite{kipf2017} and GraphSAGE~\cite{hamilton2017} aggregate neighbour information with learned weights. On tabular problems of moderate size, however, tree ensembles remain hard to beat~\cite{grinsztajn2022,shwartz2022}, and on Elliptic the evidence for a graph advantage has weakened under stricter evaluation. In a recent re-evaluation, Maganti~\cite{maganti2026} reports that a Random Forest on raw features outperforms all tested graph networks under a strictly inductive protocol, attributes a large part of previously reported graph gains to training-time exposure to test-period adjacency, and finds that randomly rewired graphs can outperform the real one under temporal shift. Malik~\cite{malik2026} compares transaction-level and actor-level review queues on Elliptic++ under a causal temporal protocol and observes, as we do independently, that actor features are constant across repeated appearances of an address.

Our study is complementary to both. We do not train a graph network; we ask how much a tree ensemble gains from \emph{precomputed} neighbourhood statistics, which amounts to one mean-aggregation step without learned weights, and we ask this separately for the transaction graph and the address graph. We also test whether the gain survives when the neighbourhood is restricted to entities that already existed, which addresses the adjacency-exposure concern raised in~\cite{maganti2026}.

\subsection{Evaluation over time, drift, and leakage}
A classifier deployed for security monitoring is trained on the past and applied to the future. Evaluations that ignore this ordering overstate performance; Pendlebury \emph{et al.}~\cite{pendlebury2019} showed this for malware classification and formulated constraints for time-aware evaluation, and Arp \emph{et al.}~\cite{arp2022} list temporal snooping, inappropriate baselines and base-rate neglect among the recurring pitfalls of machine learning in security. Kapoor and Narayanan~\cite{kapoor2023} document how leakage undermines reproducibility across scientific fields that use machine learning. Changes in the data-generating process after training are studied as concept drift~\cite{gama2014} or, more generally, data-set shift~\cite{moreno2012}; evaluating a model that is updated as labels arrive is known as prequential evaluation~\cite{gama2013}, and fraud-detection work has stressed that labels arrive with delay~\cite{dalpozzolo2018}.

For Bitcoin specifically, continual-learning approaches to anti-money-laundering have been reviewed by Deprez \emph{et al.}~\cite{deprez2025}, dynamic graph models have been evaluated around the shutdown event~\cite{kulkarni2025}, and Bhatt and Sharma~\cite{bhatt2026} argue from rolling evaluations that the dominant failure under non-stationarity is the miscalibration of decision thresholds, not the loss of predictive accuracy. Our results on the same event point the other way for a model with a fixed threshold (Section~\ref{sec:res-drift}), and we discuss why the two observations are compatible.

\subsection{Flow-based network anomaly detection}
Network intrusion detection has a long-standing gap between laboratory results and operational use~\cite{sommer2010,apruzzese2023}, and a recurring criticism concerns the data sets~\cite{ring2019}. UGR'16~\cite{macia2018} was built to address part of this: it contains anonymised NetFlow traces from a tier-3 Internet service provider, with a calibration period of real background traffic and a test period into which labelled attack traffic was injected. Its authors analyse it with multivariate statistical monitoring over counters derived from the flows~\cite{camacho2019}; Mag\'an-Carri\'on \emph{et al.}~\cite{magan2020} use it to propose a methodology for comparing learning-based detectors; and Macko \emph{et al.}~\cite{macko2023} combine a supervised filter with an autoencoder on its flows. Its botnet class is labelled Neris, a botnet that is also documented in the CTU-13 collection of botnet captures~\cite{garcia2014}. We use a public sample of the first days of the test period only, which limits what our network layer can show; Section~\ref{sec:threats} is explicit about this.

\subsection{Unsupervised detection}
Because labels are scarce in both domains, unsupervised detectors are an attractive default~\cite{chandola2009,akoglu2015}. We include Isolation Forest~\cite{liu2008} as the unsupervised reference in every layer. Its behaviour on the transaction data turns out to be informative precisely because it fails.
